\documentclass[]{extarticle}
\usepackage{stackengine,xcolor}

\newcommand\shadowfy[1]{\expandafter\shadowfypars#1\par\relax\relax}
\long\def\shadowfypars#1\par#2\relax{%
  \ifx#1\relax\else
    \shadowfywords#1 \relax\relax%
  \fi%
  \ifx\relax#2\else\par\shadowfypars#2\relax\fi%
}
\def\shadowfywords#1 #2\relax{%
  \ifx#1\relax\else
    \shadowfyletters#1\relax\relax%
  \fi%
  \ifx\relax#2\else\ \shadowfywords#2\relax\fi%
}
\def\shadowfyletters#1#2\relax{%
  \shadow{#1}%
  \ifx\relax#2\else\shadowfyletters#2\relax\fi}

\newlength\shadowHoffset
\newlength\shadowVoffset
\setlength\shadowHoffset{.2pt}
\setlength\shadowVoffset{.1pt}
\def\primarycolor{white}
\def\secondarycolor{black}

\def\shadow#1{\setstackgap{L}{0pt}\def\stacktype{L}%
\def\useanchorwidth{T}% CAN BE COMMENTEDD FOR MORE INTERLETTER SPACE.
\Longstack{%
\raisebox{0pt}{\textcolor{\primarycolor}{#1}} 
\kern.7\shadowHoffset\raisebox{.7\shadowVoffset}{\textcolor{\secondarycolor}{#1}}
\kern-.7\shadowHoffset\raisebox{.7\shadowVoffset}{\textcolor{\secondarycolor}{#1}}
\kern\shadowHoffset\raisebox{0pt}{\textcolor{\secondarycolor}{#1}}
\kern-\shadowHoffset\raisebox{0pt}{\textcolor{\secondarycolor}{#1}}
\kern.7\shadowHoffset\raisebox{-.7\shadowVoffset}{\textcolor{\secondarycolor}{#1}}
\kern-.7\shadowHoffset\raisebox{-.7\shadowVoffset}{\textcolor{\secondarycolor}{#1}}
\kern0pt\raisebox{\shadowVoffset}{\textcolor{\secondarycolor}{#1}}
\kern0pt\raisebox{-\shadowVoffset}{\textcolor{\secondarycolor}{#1}}%
}}
\usepackage[paperheight=15in,paperwidth=15in,margin=0.5in]{geometry}
\usepackage{tikz}
\usepackage{fontspec}
\usepackage{polyglossia}
\setdefaultlanguage{thai}
\setotherlanguages{english,russian}
%\setdefaultlanguage{russian}
%\setmainfont[Ligatures=TeX]{Times New Roman}
\newfontfamily\cyrillicfont{Tahoma Bold}[Script=Cyrillic]
\newfontfamily\thaifont{Sathu}[Script=Thai,Scale=MatchLowercase,WordSpace=1.25,Mapping=tex-text]
\usepackage{xunicode}
\usepackage{xltxtra}
\pagestyle{empty}
\begin{document}
\begin{tikzpicture}[remember picture,overlay]
% draw image
\node[inner sep=0] at (current page.center)
{\includegraphics[width=\paperwidth,height=\paperheight]{thaicard.png}};
\end{tikzpicture}
\def\primarycolor{orange!70}
\def\secondarycolor{black}
\resizebox*{\textwidth}{!}{\shadowfy{คำศัพท์ วลี และประโยคต่างๆที่เกี่ยวกับ “รูปลักษณ์”}}
\hspace{0pt}
\vfill
\begin{center}
\resizebox*{\textwidth}{!}{\textrussian{\textbf{телослож\textcolor{red}{\'е}ние}}}
\vskip0.7in
\resizebox*{\textwidth}{!}{[ประเ-ภทข-\textcolor{red}{องร่า}-งกาย]}
\vskip0.7in
\resizebox*{\textwidth}{!}{\textbf{(ประเภทของร่า)งกาย}}
\end{center}
\vfill
\hspace{0pt}
\pagebreak
\end{document}
